\chapter{Kết quả chi tiết theo từng hạt giống}
\label{AppendixSeedResults}

\section{Phạm vi trình bày}

Phụ lục này công bố các kết quả riêng của ba hạt giống ngẫu nhiên 42, 123 và 456 nhằm bổ trợ cho các giá trị trung bình và độ lệch chuẩn được trình bày trong Chương \ref{Chapter4}. Các bảng giữ nguyên giá trị tại từng lần huấn luyện, không lựa chọn hạt giống có kết quả thuận lợi và không dùng các giá trị này như những quan sát độc lập ở cấp mẫu. Vì chỉ có ba lần huấn luyện, sự thay đổi giữa các hạt giống được sử dụng chủ yếu để đánh giá độ nhạy của quy trình huấn luyện và hỗ trợ việc diễn giải khoảng tin cậy, thay vì khẳng định phân phối hiệu năng của mô hình.

Các kết quả được chọn gồm phân loại khi sử dụng toàn bộ dữ liệu, phân loại trong trường hợp mẫu học hạn chế, nghiên cứu thành phần leave-one-out và OOD hậu xử lý. Đây là các nội dung trong Chương \ref{Chapter4} có đủ kết quả cho cả ba hạt giống và có thể đối chiếu trực tiếp với các bảng tổng hợp. Kết quả zero-shot, phép đo hiệu quả suy luận và hai trường hợp Integrated Gradients không được lặp lại dưới dạng bảng ba hạt giống vì các nội dung này hiện chỉ được thực hiện một lần hoặc được dùng như minh họa định tính; cách trình bày này tránh tạo ấn tượng rằng các phép đo đơn lẻ đã được lặp lại.

\section{Phân loại khi sử dụng toàn bộ dữ liệu huấn luyện}

Bảng \ref{tab:appendix-full-shot-seeds-btxrd} và Bảng \ref{tab:appendix-full-shot-seeds-ctch} trình bày năm độ đo phân loại của từng mô hình tại từng hạt giống. Việc giữ đồng thời Accuracy, Balanced Accuracy, Macro-F1, Macro-AUROC và Macro-AUPRC cho phép kiểm tra liệu thứ hạng trung bình trong Chương \ref{Chapter4} có xuất hiện nhất quán hay phụ thuộc vào một lần khởi tạo cụ thể.

\begingroup
\scriptsize
\setlength{\tabcolsep}{3pt}
\renewcommand{\arraystretch}{1.08}
\begin{longtable}{L{3.2cm}C{1.15cm}*{5}{C{1.55cm}}}
\caption{Kết quả full-shot của từng hạt giống trên BTXRD; hàng XBone-Net được tô xám để dễ đối chiếu.}\label{tab:appendix-full-shot-seeds-btxrd}\\
\toprule
\textbf{Mô hình} & \textbf{Hạt giống} & \textbf{Accuracy} & \textbf{BAcc} & \textbf{Macro-F1} & \textbf{Macro-AUROC} & \textbf{Macro-AUPRC} \\
\midrule
\endfirsthead
\multicolumn{7}{c}{\tablename\ \thetable\ (tiếp theo)}\\
\toprule
\textbf{Mô hình} & \textbf{Hạt giống} & \textbf{Accuracy} & \textbf{BAcc} & \textbf{Macro-F1} & \textbf{Macro-AUROC} & \textbf{Macro-AUPRC} \\
\midrule
\endhead
\midrule \multicolumn{7}{r}{Còn tiếp ở trang sau}\\
\endfoot
\bottomrule
\endlastfoot
BioMedCLIP & 42 & 0{,}8320 & 0{,}5769 & 0{,}5870 & 0{,}8679 & 0{,}6395 \\
BioMedCLIP & 123 & 0{,}8280 & 0{,}6197 & 0{,}5993 & 0{,}8918 & 0{,}6542 \\
BioMedCLIP & 456 & 0{,}8333 & 0{,}6219 & 0{,}6118 & 0{,}8996 & 0{,}6595 \\
CLIP & 42 & 0{,}8200 & 0{,}5949 & 0{,}5722 & 0{,}9114 & 0{,}6342 \\
CLIP & 123 & 0{,}7720 & 0{,}5802 & 0{,}5574 & 0{,}8706 & 0{,}6192 \\
CLIP & 456 & 0{,}8187 & 0{,}6121 & 0{,}6064 & 0{,}9099 & 0{,}6219 \\
DenseNet-121 & 42 & 0{,}6320 & 0{,}3644 & 0{,}3688 & 0{,}7916 & 0{,}3852 \\
DenseNet-121 & 123 & 0{,}6333 & 0{,}3862 & 0{,}3904 & 0{,}7596 & 0{,}3790 \\
DenseNet-121 & 456 & 0{,}6307 & 0{,}3575 & 0{,}3562 & 0{,}7639 & 0{,}3851 \\
MedCLIP & 42 & 0{,}8307 & 0{,}5860 & 0{,}5931 & 0{,}8221 & 0{,}6089 \\
MedCLIP & 123 & 0{,}8040 & 0{,}6204 & 0{,}6014 & 0{,}8736 & 0{,}6074 \\
MedCLIP & 456 & 0{,}8267 & 0{,}6310 & 0{,}6217 & 0{,}8838 & 0{,}6520 \\
PubMedCLIP & 42 & 0{,}7573 & 0{,}5716 & 0{,}5067 & 0{,}8856 & 0{,}5874 \\
PubMedCLIP & 123 & 0{,}7947 & 0{,}6163 & 0{,}5946 & 0{,}8970 & 0{,}6233 \\
PubMedCLIP & 456 & 0{,}8133 & 0{,}6388 & 0{,}6155 & 0{,}9317 & 0{,}6341 \\
ResNet-50 & 42 & 0{,}6067 & 0{,}3031 & 0{,}3174 & 0{,}6947 & 0{,}3279 \\
ResNet-50 & 123 & 0{,}6213 & 0{,}3253 & 0{,}3364 & 0{,}7447 & 0{,}3401 \\
ResNet-50 & 456 & 0{,}5933 & 0{,}3278 & 0{,}3333 & 0{,}7313 & 0{,}3333 \\
LoRA-BiomedCLIP & 42 & 0{,}8173 & 0{,}6402 & 0{,}6226 & 0{,}8859 & 0{,}6545 \\
LoRA-BiomedCLIP & 123 & 0{,}8107 & 0{,}6421 & 0{,}6136 & 0{,}8984 & 0{,}6127 \\
LoRA-BiomedCLIP & 456 & 0{,}8120 & 0{,}6834 & 0{,}6312 & 0{,}9463 & 0{,}6357 \\
LoRA-PubMedCLIP & 42 & 0{,}8280 & 0{,}6227 & 0{,}6177 & 0{,}8975 & 0{,}6337 \\
LoRA-PubMedCLIP & 123 & 0{,}8373 & 0{,}6594 & 0{,}6569 & 0{,}9100 & 0{,}6979 \\
LoRA-PubMedCLIP & 456 & 0{,}8387 & 0{,}6287 & 0{,}6237 & 0{,}9043 & 0{,}6460 \\
\rowcolor{gray!12}XBone-Net & 42 & 0{,}8227 & 0{,}6029 & 0{,}6037 & 0{,}9391 & 0{,}6520 \\
\rowcolor{gray!12}XBone-Net & 123 & 0{,}8253 & 0{,}5797 & 0{,}5972 & 0{,}8899 & 0{,}6099 \\
\rowcolor{gray!12}XBone-Net & 456 & 0{,}8293 & 0{,}5922 & 0{,}6023 & 0{,}9438 & 0{,}6285 \\
\end{longtable}
\endgroup

\begingroup
\scriptsize
\setlength{\tabcolsep}{3pt}
\renewcommand{\arraystretch}{1.08}
\begin{longtable}{L{3.2cm}C{1.15cm}*{5}{C{1.55cm}}}
\caption{Kết quả full-shot của từng hạt giống trên CTCH; hàng XBone-Net được tô xám để dễ đối chiếu.}\label{tab:appendix-full-shot-seeds-ctch}\\
\toprule
\textbf{Mô hình} & \textbf{Hạt giống} & \textbf{Accuracy} & \textbf{BAcc} & \textbf{Macro-F1} & \textbf{Macro-AUROC} & \textbf{Macro-AUPRC} \\
\midrule
\endfirsthead
\multicolumn{7}{c}{\tablename\ \thetable\ (tiếp theo)}\\
\toprule
\textbf{Mô hình} & \textbf{Hạt giống} & \textbf{Accuracy} & \textbf{BAcc} & \textbf{Macro-F1} & \textbf{Macro-AUROC} & \textbf{Macro-AUPRC} \\
\midrule
\endhead
\midrule \multicolumn{7}{r}{Còn tiếp ở trang sau}\\
\endfoot
\bottomrule
\endlastfoot
BioMedCLIP & 42 & 0{,}5904 & 0{,}6203 & 0{,}5520 & 0{,}9559 & 0{,}6248 \\
BioMedCLIP & 123 & 0{,}5755 & 0{,}6107 & 0{,}5472 & 0{,}9555 & 0{,}6188 \\
BioMedCLIP & 456 & 0{,}6248 & 0{,}5758 & 0{,}5628 & 0{,}9057 & 0{,}6071 \\
CLIP & 42 & 0{,}5725 & 0{,}6579 & 0{,}5727 & 0{,}9547 & 0{,}6270 \\
CLIP & 123 & 0{,}5785 & 0{,}6209 & 0{,}5547 & 0{,}9478 & 0{,}6196 \\
CLIP & 456 & 0{,}6114 & 0{,}5862 & 0{,}5624 & 0{,}8845 & 0{,}5542 \\
DenseNet-121 & 42 & 0{,}4813 & 0{,}5348 & 0{,}4610 & 0{,}9404 & 0{,}4917 \\
DenseNet-121 & 123 & 0{,}4873 & 0{,}4435 & 0{,}4134 & 0{,}9323 & 0{,}4667 \\
DenseNet-121 & 456 & 0{,}4768 & 0{,}4999 & 0{,}4356 & 0{,}9325 & 0{,}4703 \\
MedCLIP & 42 & 0{,}5590 & 0{,}5811 & 0{,}5265 & 0{,}9354 & 0{,}5973 \\
MedCLIP & 123 & 0{,}5979 & 0{,}5503 & 0{,}5285 & 0{,}8835 & 0{,}5628 \\
MedCLIP & 456 & 0{,}5202 & 0{,}5838 & 0{,}4902 & 0{,}9456 & 0{,}5916 \\
PubMedCLIP & 42 & 0{,}5306 & 0{,}6108 & 0{,}5313 & 0{,}9350 & 0{,}6073 \\
PubMedCLIP & 123 & 0{,}5486 & 0{,}4912 & 0{,}4902 & 0{,}9038 & 0{,}5683 \\
PubMedCLIP & 456 & 0{,}5770 & 0{,}5747 & 0{,}5395 & 0{,}9448 & 0{,}5907 \\
ResNet-50 & 42 & 0{,}4918 & 0{,}4238 & 0{,}4131 & 0{,}9130 & 0{,}4489 \\
ResNet-50 & 123 & 0{,}4948 & 0{,}4027 & 0{,}4029 & 0{,}9118 & 0{,}4452 \\
ResNet-50 & 456 & 0{,}4918 & 0{,}4114 & 0{,}4076 & 0{,}9188 & 0{,}4536 \\
LoRA-BiomedCLIP & 42 & 0{,}6203 & 0{,}6113 & 0{,}5762 & 0{,}9095 & 0{,}5829 \\
LoRA-BiomedCLIP & 123 & 0{,}5830 & 0{,}6324 & 0{,}5685 & 0{,}9516 & 0{,}6235 \\
LoRA-BiomedCLIP & 456 & 0{,}5845 & 0{,}5839 & 0{,}5349 & 0{,}9389 & 0{,}6071 \\
LoRA-PubMedCLIP & 42 & 0{,}5665 & 0{,}5718 & 0{,}5314 & 0{,}9600 & 0{,}6031 \\
LoRA-PubMedCLIP & 123 & 0{,}6114 & 0{,}5769 & 0{,}5631 & 0{,}9253 & 0{,}6167 \\
LoRA-PubMedCLIP & 456 & 0{,}6248 & 0{,}5610 & 0{,}5671 & 0{,}9246 & 0{,}6062 \\
\rowcolor{gray!12}XBone-Net & 42 & 0{,}5590 & 0{,}5995 & 0{,}5485 & 0{,}9601 & 0{,}6260 \\
\rowcolor{gray!12}XBone-Net & 123 & 0{,}6054 & 0{,}5334 & 0{,}5629 & 0{,}9190 & 0{,}6022 \\
\rowcolor{gray!12}XBone-Net & 456 & 0{,}5531 & 0{,}6203 & 0{,}5559 & 0{,}9628 & 0{,}6284 \\
\end{longtable}
\endgroup


Hình \ref{fig:appendix-full-shot-f1-seeds} tập trung vào XBone-Net và hai mô hình cơ sở PEFT gần nhất để minh họa biến thiên Macro-F1 mà không lặp lại toàn bộ thông tin trong hai bảng. Trên BTXRD, đường của XBone-Net thay đổi ít hơn hai mô hình cơ sở nhưng nằm ở mức thấp hơn; trên CTCH, ba đường giao nhau giữa các hạt giống, phù hợp với nhận định rằng chênh lệch Macro-F1 giữa các mô hình ngôn ngữ trực quan chưa ổn định trong phạm vi ba lần huấn luyện.

\begin{figure}[H]
\centering
\includegraphics[width=\textwidth]{generated/appendix_seed_results/full_shot_macro_f1_by_seed.png}
\caption{Macro-F1 full-shot của XBone-Net và hai mô hình cơ sở PEFT tại từng hạt giống trên BTXRD và CTCH}
\label{fig:appendix-full-shot-f1-seeds}
\end{figure}

\section{Phân loại trong trường hợp mẫu học hạn chế}

Bảng \ref{tab:appendix-few-shot-seeds-btxrd} và Bảng \ref{tab:appendix-few-shot-seeds-ctch} báo cáo kết quả tại từng hạt giống cho các mức ngân sách tối đa 1, 10 và 20 mẫu trên mỗi lớp. Khi số mẫu huấn luyện của một lớp thấp hơn ngân sách, toàn bộ mẫu hiện có được giữ lại và không lấy lặp; vì vậy mức 20 mẫu của CTCH sử dụng 19 mẫu cho lớp gãy xương thuyền và 16 mẫu cho lớp gãy xương vùng chậu. Các giá trị này bổ sung cho biểu đồ trung bình trong Chương \ref{Chapter4}, đồng thời cho phép nhận diện trường hợp chênh lệch trung bình nhỏ hơn mức dao động giữa các lần huấn luyện.

\begin{table}[htbp]
\centering
\footnotesize
\renewcommand{\arraystretch}{1.08}
\resizebox{\textwidth}{!}{%
\begin{tabular}{C{1.0cm}L{2.35cm}C{1.0cm}*{5}{C{1.45cm}}}
\toprule
\textbf{Thiết lập} & \textbf{Mô hình} & \textbf{Hạt giống} & \textbf{Accuracy} & \textbf{BAcc} & \textbf{Macro-F1} & \textbf{Macro-AUROC} & \textbf{Macro-AUPRC} \\
\midrule
1 mẫu & LoRA-BiomedCLIP & 42 & 0{,}2880 & 0{,}1666 & 0{,}1413 & 0{,}6239 & 0{,}1747 \\
1 mẫu & LoRA-BiomedCLIP & 123 & 0{,}3280 & 0{,}2601 & 0{,}2098 & 0{,}6756 & 0{,}2430 \\
1 mẫu & LoRA-BiomedCLIP & 456 & 0{,}2187 & 0{,}2130 & 0{,}1587 & 0{,}5960 & 0{,}2077 \\
1 mẫu & LoRA-PubMedCLIP & 42 & 0{,}5027 & 0{,}2905 & 0{,}2506 & 0{,}6672 & 0{,}2610 \\
1 mẫu & LoRA-PubMedCLIP & 123 & 0{,}2840 & 0{,}2566 & 0{,}1643 & 0{,}6649 & 0{,}2799 \\
1 mẫu & LoRA-PubMedCLIP & 456 & 0{,}1627 & 0{,}2050 & 0{,}1304 & 0{,}6213 & 0{,}2095 \\
\rowcolor{gray!12}1 mẫu & XBone-Net & 42 & 0{,}2053 & 0{,}1547 & 0{,}1345 & 0{,}5958 & 0{,}1608 \\
\rowcolor{gray!12}1 mẫu & XBone-Net & 123 & 0{,}2720 & 0{,}2499 & 0{,}1859 & 0{,}6450 & 0{,}2347 \\
\rowcolor{gray!12}1 mẫu & XBone-Net & 456 & 0{,}2667 & 0{,}2302 & 0{,}1655 & 0{,}6413 & 0{,}2150 \\
\bottomrule
\end{tabular}%
}
\caption{Kết quả mẫu học hạn chế của từng hạt giống trên BTXRD ở thiết lập 1 mẫu}
\label{tab:appendix-few-shot-seeds-btxrd}
\end{table}

\begin{table}[htbp]
\ContinuedFloat
\centering
\footnotesize
\renewcommand{\arraystretch}{1.08}
\resizebox{\textwidth}{!}{%
\begin{tabular}{C{1.0cm}L{2.35cm}C{1.0cm}*{5}{C{1.45cm}}}
\toprule
\textbf{Thiết lập} & \textbf{Mô hình} & \textbf{Hạt giống} & \textbf{Accuracy} & \textbf{BAcc} & \textbf{Macro-F1} & \textbf{Macro-AUROC} & \textbf{Macro-AUPRC} \\
\midrule
10 mẫu & LoRA-BiomedCLIP & 42 & 0{,}6373 & 0{,}4487 & 0{,}4074 & 0{,}8503 & 0{,}4289 \\
10 mẫu & LoRA-BiomedCLIP & 123 & 0{,}7133 & 0{,}4759 & 0{,}4308 & 0{,}8764 & 0{,}4634 \\
10 mẫu & LoRA-BiomedCLIP & 456 & 0{,}6800 & 0{,}4766 & 0{,}4251 & 0{,}8528 & 0{,}4480 \\
10 mẫu & LoRA-PubMedCLIP & 42 & 0{,}6160 & 0{,}4535 & 0{,}4010 & 0{,}8366 & 0{,}4533 \\
10 mẫu & LoRA-PubMedCLIP & 123 & 0{,}7333 & 0{,}4460 & 0{,}4159 & 0{,}8738 & 0{,}4814 \\
10 mẫu & LoRA-PubMedCLIP & 456 & 0{,}6427 & 0{,}4328 & 0{,}4324 & 0{,}8483 & 0{,}4471 \\
\rowcolor{gray!12}10 mẫu & XBone-Net & 42 & 0{,}6547 & 0{,}4894 & 0{,}4255 & 0{,}8599 & 0{,}4463 \\
\rowcolor{gray!12}10 mẫu & XBone-Net & 123 & 0{,}7067 & 0{,}4429 & 0{,}4066 & 0{,}8520 & 0{,}4276 \\
\rowcolor{gray!12}10 mẫu & XBone-Net & 456 & 0{,}7053 & 0{,}4989 & 0{,}4585 & 0{,}8612 & 0{,}4918 \\
\bottomrule
\end{tabular}%
}
\caption[]{Kết quả mẫu học hạn chế của từng hạt giống trên BTXRD ở thiết lập 10 mẫu}
\end{table}

\begin{table}[htbp]
\ContinuedFloat
\centering
\footnotesize
\renewcommand{\arraystretch}{1.08}
\resizebox{\textwidth}{!}{%
\begin{tabular}{C{1.0cm}L{2.35cm}C{1.0cm}*{5}{C{1.45cm}}}
\toprule
\textbf{Thiết lập} & \textbf{Mô hình} & \textbf{Hạt giống} & \textbf{Accuracy} & \textbf{BAcc} & \textbf{Macro-F1} & \textbf{Macro-AUROC} & \textbf{Macro-AUPRC} \\
\midrule
20 mẫu & LoRA-BiomedCLIP & 42 & 0{,}7347 & 0{,}5774 & 0{,}4866 & 0{,}9268 & 0{,}5491 \\
20 mẫu & LoRA-BiomedCLIP & 123 & 0{,}7360 & 0{,}5353 & 0{,}4838 & 0{,}9365 & 0{,}5383 \\
20 mẫu & LoRA-BiomedCLIP & 456 & 0{,}7227 & 0{,}5350 & 0{,}4925 & 0{,}9181 & 0{,}5388 \\
20 mẫu & LoRA-PubMedCLIP & 42 & 0{,}7293 & 0{,}5608 & 0{,}5062 & 0{,}9119 & 0{,}5608 \\
20 mẫu & LoRA-PubMedCLIP & 123 & 0{,}7320 & 0{,}5837 & 0{,}5073 & 0{,}9357 & 0{,}5592 \\
20 mẫu & LoRA-PubMedCLIP & 456 & 0{,}7493 & 0{,}5911 & 0{,}5261 & 0{,}9139 & 0{,}5484 \\
\rowcolor{gray!12}20 mẫu & XBone-Net & 42 & 0{,}7347 & 0{,}5794 & 0{,}5076 & 0{,}8991 & 0{,}5410 \\
\rowcolor{gray!12}20 mẫu & XBone-Net & 123 & 0{,}7360 & 0{,}5806 & 0{,}4987 & 0{,}9192 & 0{,}5403 \\
\rowcolor{gray!12}20 mẫu & XBone-Net & 456 & 0{,}7640 & 0{,}6246 & 0{,}5545 & 0{,}9307 & 0{,}6091 \\
\bottomrule
\end{tabular}%
}
\caption[]{Kết quả mẫu học hạn chế của từng hạt giống trên BTXRD ở thiết lập 20 mẫu}
\end{table}

\begin{table}[htbp]
\centering
\footnotesize
\renewcommand{\arraystretch}{1.08}
\resizebox{\textwidth}{!}{%
\begin{tabular}{C{1.0cm}L{2.35cm}C{1.0cm}*{5}{C{1.45cm}}}
\toprule
\textbf{Thiết lập} & \textbf{Mô hình} & \textbf{Hạt giống} & \textbf{Accuracy} & \textbf{BAcc} & \textbf{Macro-F1} & \textbf{Macro-AUROC} & \textbf{Macro-AUPRC} \\
\midrule
1 mẫu & LoRA-BiomedCLIP & 42 & 0{,}2093 & 0{,}2674 & 0{,}2110 & 0{,}8286 & 0{,}2244 \\
1 mẫu & LoRA-BiomedCLIP & 123 & 0{,}2257 & 0{,}3002 & 0{,}2067 & 0{,}8411 & 0{,}2328 \\
1 mẫu & LoRA-BiomedCLIP & 456 & 0{,}2033 & 0{,}2814 & 0{,}2044 & 0{,}7997 & 0{,}2500 \\
1 mẫu & LoRA-PubMedCLIP & 42 & 0{,}1480 & 0{,}2108 & 0{,}1435 & 0{,}7689 & 0{,}1781 \\
1 mẫu & LoRA-PubMedCLIP & 123 & 0{,}2451 & 0{,}3605 & 0{,}2303 & 0{,}8032 & 0{,}2294 \\
1 mẫu & LoRA-PubMedCLIP & 456 & 0{,}1794 & 0{,}2839 & 0{,}1890 & 0{,}7991 & 0{,}2044 \\
\rowcolor{gray!12}1 mẫu & XBone-Net & 42 & 0{,}2272 & 0{,}2825 & 0{,}2075 & 0{,}8662 & 0{,}2781 \\
\rowcolor{gray!12}1 mẫu & XBone-Net & 123 & 0{,}2541 & 0{,}4090 & 0{,}2467 & 0{,}8718 & 0{,}3041 \\
\rowcolor{gray!12}1 mẫu & XBone-Net & 456 & 0{,}2212 & 0{,}3137 & 0{,}2165 & 0{,}8556 & 0{,}2823 \\
\bottomrule
\end{tabular}%
}
\caption{Kết quả mẫu học hạn chế của từng hạt giống trên CTCH ở thiết lập 1 mẫu}
\label{tab:appendix-few-shot-seeds-ctch}
\end{table}

\begin{table}[htbp]
\ContinuedFloat
\centering
\footnotesize
\renewcommand{\arraystretch}{1.08}
\resizebox{\textwidth}{!}{%
\begin{tabular}{C{1.0cm}L{2.35cm}C{1.0cm}*{5}{C{1.45cm}}}
\toprule
\textbf{Thiết lập} & \textbf{Mô hình} & \textbf{Hạt giống} & \textbf{Accuracy} & \textbf{BAcc} & \textbf{Macro-F1} & \textbf{Macro-AUROC} & \textbf{Macro-AUPRC} \\
\midrule
10 mẫu & LoRA-BiomedCLIP & 42 & 0{,}4141 & 0{,}5817 & 0{,}4384 & 0{,}9399 & 0{,}4741 \\
10 mẫu & LoRA-BiomedCLIP & 123 & 0{,}3662 & 0{,}5270 & 0{,}4101 & 0{,}9272 & 0{,}4345 \\
10 mẫu & LoRA-BiomedCLIP & 456 & 0{,}4305 & 0{,}5588 & 0{,}4389 & 0{,}9387 & 0{,}4855 \\
10 mẫu & LoRA-PubMedCLIP & 42 & 0{,}3871 & 0{,}5763 & 0{,}4281 & 0{,}9452 & 0{,}4784 \\
10 mẫu & LoRA-PubMedCLIP & 123 & 0{,}3632 & 0{,}5491 & 0{,}4055 & 0{,}9323 & 0{,}4732 \\
10 mẫu & LoRA-PubMedCLIP & 456 & 0{,}3931 & 0{,}5633 & 0{,}4220 & 0{,}9404 & 0{,}4837 \\
\rowcolor{gray!12}10 mẫu & XBone-Net & 42 & 0{,}4111 & 0{,}5928 & 0{,}4528 & 0{,}9414 & 0{,}4627 \\
\rowcolor{gray!12}10 mẫu & XBone-Net & 123 & 0{,}3348 & 0{,}5122 & 0{,}3917 & 0{,}9245 & 0{,}4274 \\
\rowcolor{gray!12}10 mẫu & XBone-Net & 456 & 0{,}4320 & 0{,}5721 & 0{,}4542 & 0{,}9445 & 0{,}4720 \\
\bottomrule
\end{tabular}%
}
\caption[]{Kết quả mẫu học hạn chế của từng hạt giống trên CTCH ở thiết lập 10 mẫu}
\end{table}

\begin{table}[htbp]
\ContinuedFloat
\centering
\footnotesize
\renewcommand{\arraystretch}{1.08}
\resizebox{\textwidth}{!}{%
\begin{tabular}{C{1.0cm}L{2.35cm}C{1.0cm}*{5}{C{1.45cm}}}
\toprule
\textbf{Thiết lập} & \textbf{Mô hình} & \textbf{Hạt giống} & \textbf{Accuracy} & \textbf{BAcc} & \textbf{Macro-F1} & \textbf{Macro-AUROC} & \textbf{Macro-AUPRC} \\
\midrule
20 mẫu & LoRA-BiomedCLIP & 42 & 0{,}4694 & 0{,}6409 & 0{,}5033 & 0{,}9527 & 0{,}5364 \\
20 mẫu & LoRA-BiomedCLIP & 123 & 0{,}4170 & 0{,}5883 & 0{,}4597 & 0{,}9255 & 0{,}4756 \\
20 mẫu & LoRA-BiomedCLIP & 456 & 0{,}4320 & 0{,}6117 & 0{,}4777 & 0{,}9466 & 0{,}5142 \\
20 mẫu & LoRA-PubMedCLIP & 42 & 0{,}4694 & 0{,}6076 & 0{,}4899 & 0{,}9495 & 0{,}5319 \\
20 mẫu & LoRA-PubMedCLIP & 123 & 0{,}4305 & 0{,}5932 & 0{,}4430 & 0{,}9493 & 0{,}5050 \\
20 mẫu & LoRA-PubMedCLIP & 456 & 0{,}4529 & 0{,}5626 & 0{,}4618 & 0{,}9513 & 0{,}5420 \\
\rowcolor{gray!12}20 mẫu & XBone-Net & 42 & 0{,}4559 & 0{,}5725 & 0{,}4756 & 0{,}9374 & 0{,}5223 \\
\rowcolor{gray!12}20 mẫu & XBone-Net & 123 & 0{,}4200 & 0{,}6052 & 0{,}4659 & 0{,}9494 & 0{,}5067 \\
\rowcolor{gray!12}20 mẫu & XBone-Net & 456 & 0{,}4320 & 0{,}5789 & 0{,}4637 & 0{,}9503 & 0{,}5362 \\
\bottomrule
\end{tabular}%
}
\caption[]{Kết quả mẫu học hạn chế của từng hạt giống trên CTCH ở thiết lập 20 mẫu}
\end{table}


\section{Nghiên cứu thành phần leave-one-out}

Bảng \ref{tab:appendix-ablation-seeds} trình bày kết quả của XBone-Net và năm biến thể leave-one-out tại từng hạt giống trên CTCH. Bảng này được dùng để đối chiếu với các forest plot trong Chương \ref{Chapter4}; một biến thể có giá trị trung bình cao hơn nhưng thay đổi dấu chênh lệch giữa các hạt giống không được xem là bằng chứng về ưu thế ổn định.

\begin{table}[htbp]
\centering
\footnotesize
\renewcommand{\arraystretch}{1.08}
\resizebox{\textwidth}{!}{%
\begin{tabular}{L{2.55cm}C{1.0cm}*{6}{C{1.45cm}}}
\toprule
\textbf{Cấu hình} & \textbf{Hạt giống} & \textbf{Accuracy} & \textbf{BAcc} & \textbf{Macro-F1} & \textbf{Macro-AUROC} & \textbf{Macro-AUPRC} & \textbf{ECE} \\
\midrule
\rowcolor{gray!12}XBone-Net & 42 & 0{,}5590 & 0{,}5995 & 0{,}5485 & 0{,}9601 & 0{,}6260 & 0{,}0857 \\
\rowcolor{gray!12}XBone-Net & 123 & 0{,}6054 & 0{,}5334 & 0{,}5629 & 0{,}9190 & 0{,}6022 & 0{,}0856 \\
\rowcolor{gray!12}XBone-Net & 456 & 0{,}5531 & 0{,}6203 & 0{,}5559 & 0{,}9628 & 0{,}6284 & 0{,}1382 \\
Không dùng ảnh độ phân giải cao & 42 & 0{,}6293 & 0{,}5454 & 0{,}5512 & 0{,}8943 & 0{,}6022 & 0{,}1502 \\
Không dùng ảnh độ phân giải cao & 123 & 0{,}6039 & 0{,}5608 & 0{,}5704 & 0{,}9116 & 0{,}5907 & 0{,}2462 \\
Không dùng ảnh độ phân giải cao & 456 & 0{,}5889 & 0{,}5651 & 0{,}5442 & 0{,}9385 & 0{,}6041 & 0{,}1265 \\
Gộp trung bình & 42 & 0{,}5635 & 0{,}6339 & 0{,}5732 & 0{,}9613 & 0{,}6065 & 0{,}2042 \\
Gộp trung bình & 123 & 0{,}5904 & 0{,}5520 & 0{,}5487 & 0{,}9484 & 0{,}6259 & 0{,}0918 \\
Gộp trung bình & 456 & 0{,}5830 & 0{,}6288 & 0{,}5989 & 0{,}9636 & 0{,}6612 & 0{,}1559 \\
\bottomrule
\end{tabular}%
}
\caption{Kết quả của từng hạt giống trong nghiên cứu loại bỏ từng thành phần trên CTCH (phần 1); cấu hình đầy đủ được tô xám.}
\label{tab:appendix-ablation-seeds}
\end{table}

\begin{table}[htbp]
\ContinuedFloat
\centering
\footnotesize
\renewcommand{\arraystretch}{1.08}
\resizebox{\textwidth}{!}{%
\begin{tabular}{L{2.55cm}C{1.0cm}*{6}{C{1.45cm}}}
\toprule
\textbf{Cấu hình} & \textbf{Hạt giống} & \textbf{Accuracy} & \textbf{BAcc} & \textbf{Macro-F1} & \textbf{Macro-AUROC} & \textbf{Macro-AUPRC} & \textbf{ECE} \\
\midrule
Chỉ pha 2 & 42 & 0{,}6263 & 0{,}5833 & 0{,}5950 & 0{,}9188 & 0{,}6388 & 0{,}1276 \\
Chỉ pha 2 & 123 & 0{,}6218 & 0{,}5365 & 0{,}5615 & 0{,}9317 & 0{,}6341 & 0{,}1653 \\
Chỉ pha 2 & 456 & 0{,}6039 & 0{,}5516 & 0{,}5696 & 0{,}8928 & 0{,}6073 & 0{,}2675 \\
Nối đặc trưng & 42 & 0{,}5994 & 0{,}5439 & 0{,}5357 & 0{,}9520 & 0{,}6070 & 0{,}0629 \\
Nối đặc trưng & 123 & 0{,}5949 & 0{,}5974 & 0{,}5576 & 0{,}9612 & 0{,}6124 & 0{,}0922 \\
Nối đặc trưng & 456 & 0{,}5830 & 0{,}5590 & 0{,}5344 & 0{,}9567 & 0{,}5862 & 0{,}0815 \\
Đầu tuyến tính & 42 & 0{,}6203 & 0{,}5590 & 0{,}5741 & 0{,}8984 & 0{,}5919 & 0{,}1039 \\
Đầu tuyến tính & 123 & 0{,}6323 & 0{,}5513 & 0{,}5747 & 0{,}9075 & 0{,}5985 & 0{,}0873 \\
Đầu tuyến tính & 456 & 0{,}6263 & 0{,}6234 & 0{,}6179 & 0{,}9181 & 0{,}6325 & 0{,}0733 \\
\bottomrule
\end{tabular}%
}
\caption[]{Kết quả của từng hạt giống trong nghiên cứu loại bỏ từng thành phần trên CTCH (phần 2); cấu hình đầy đủ được tô xám.}
\end{table}


\section{OOD hậu xử lý}

Bảng \ref{tab:appendix-ood-seeds} công bố kết quả theo từng hạt giống cho bốn phương pháp chấm điểm OOD trên Semantic OOD và BTXRD. Số mẫu trong phân phối và ngoài phân phối được ghi trực tiếp trong bảng vì AUPR-Out và FPR@95\%TPR chịu ảnh hưởng đáng kể bởi số lượng và tỷ lệ mẫu của từng kịch bản.

Đối với các tệp kết quả hiện có, mỗi lần đánh giá sử dụng 669 mẫu CTCH trong phân phối. Kịch bản BTXRD sử dụng đủ 750 mẫu ngoài phân phối, trong khi kịch bản Semantic OOD chỉ đọc thành công 29 trong tổng số 44 mẫu dự kiến; 15 mẫu còn lại gặp lỗi quyền truy cập tại thời điểm chạy đánh giá. Vì vậy, kết quả Semantic OOD trong Chương \ref{Chapter4} và phụ lục này phản ánh 29 mẫu đã được xử lý, không đại diện cho toàn bộ 44 mẫu ban đầu. Giới hạn này cần được cân nhắc khi diễn giải độ lệch giữa các hạt giống và đặc biệt khi so sánh AUPR-Out giữa hai kịch bản có quy mô khác nhau.

\begingroup
\scriptsize
\setlength{\tabcolsep}{3pt}
\renewcommand{\arraystretch}{1.08}
\begin{longtable}{L{1.8cm}L{2.65cm}C{0.9cm}C{0.9cm}C{0.9cm}*{3}{C{1.45cm}}}
\caption{Kết quả OOD hậu xử lý của từng hạt giống; các hàng Mahalanobis-centroid được tô xám.}\label{tab:appendix-ood-seeds}\\
\toprule
\textbf{Kịch bản} & \textbf{Phương pháp} & \textbf{Hạt giống} & \textbf{$N_{ID}$} & \textbf{$N_{OOD}$} & \textbf{AUROC-OOD} & \textbf{AUPR-Out} & \textbf{\makecell{FPR@\\95\%TPR}} \\
\midrule
\endfirsthead
\multicolumn{8}{c}{\tablename\ \thetable\ (tiếp theo)}\\
\toprule
\textbf{Kịch bản} & \textbf{Phương pháp} & \textbf{Hạt giống} & \textbf{$N_{ID}$} & \textbf{$N_{OOD}$} & \textbf{AUROC-OOD} & \textbf{AUPR-Out} & \textbf{\makecell{FPR@\\95\%TPR}} \\
\midrule
\endhead
\midrule \multicolumn{8}{r}{Còn tiếp ở trang sau}\\
\endfoot
\bottomrule
\endlastfoot
Semantic OOD & Cosine-centroid & 42 & 669 & 29 & 0{,}5596 & 0{,}0615 & 0{,}8966 \\
\rowcolor{gray!12}Semantic OOD & Mahalanobis-centroid & 42 & 669 & 29 & 0{,}5651 & 0{,}0521 & 0{,}8621 \\
Semantic OOD & kNN & 42 & 669 & 29 & 0{,}5467 & 0{,}0643 & 0{,}8966 \\
Semantic OOD & Entropy & 42 & 669 & 29 & 0{,}5572 & 0{,}0864 & 0{,}8966 \\
Semantic OOD & Cosine-centroid & 123 & 669 & 29 & 0{,}5948 & 0{,}0569 & 0{,}9310 \\
\rowcolor{gray!12}Semantic OOD & Mahalanobis-centroid & 123 & 669 & 29 & 0{,}6055 & 0{,}0570 & 0{,}9655 \\
Semantic OOD & kNN & 123 & 669 & 29 & 0{,}6034 & 0{,}0604 & 0{,}9655 \\
Semantic OOD & Entropy & 123 & 669 & 29 & 0{,}5458 & 0{,}0478 & 0{,}9655 \\
Semantic OOD & Cosine-centroid & 456 & 669 & 29 & 0{,}5674 & 0{,}0600 & 0{,}8621 \\
\rowcolor{gray!12}Semantic OOD & Mahalanobis-centroid & 456 & 669 & 29 & 0{,}5814 & 0{,}0568 & 0{,}9310 \\
Semantic OOD & kNN & 456 & 669 & 29 & 0{,}5843 & 0{,}0576 & 0{,}9310 \\
Semantic OOD & Entropy & 456 & 669 & 29 & 0{,}4964 & 0{,}0560 & 0{,}8966 \\
BTXRD & Cosine-centroid & 42 & 669 & 750 & 0{,}8052 & 0{,}7974 & 0{,}7307 \\
\rowcolor{gray!12}BTXRD & Mahalanobis-centroid & 42 & 669 & 750 & 0{,}8762 & 0{,}8591 & 0{,}6307 \\
BTXRD & kNN & 42 & 669 & 750 & 0{,}8544 & 0{,}8379 & 0{,}6573 \\
BTXRD & Entropy & 42 & 669 & 750 & 0{,}7325 & 0{,}7583 & 0{,}7453 \\
BTXRD & Cosine-centroid & 123 & 669 & 750 & 0{,}6379 & 0{,}6751 & 0{,}8160 \\
\rowcolor{gray!12}BTXRD & Mahalanobis-centroid & 123 & 669 & 750 & 0{,}6982 & 0{,}7377 & 0{,}7600 \\
BTXRD & kNN & 123 & 669 & 750 & 0{,}6654 & 0{,}6979 & 0{,}8067 \\
BTXRD & Entropy & 123 & 669 & 750 & 0{,}5931 & 0{,}6203 & 0{,}8680 \\
BTXRD & Cosine-centroid & 456 & 669 & 750 & 0{,}7471 & 0{,}7079 & 0{,}8400 \\
\rowcolor{gray!12}BTXRD & Mahalanobis-centroid & 456 & 669 & 750 & 0{,}8381 & 0{,}8058 & 0{,}7307 \\
BTXRD & kNN & 456 & 669 & 750 & 0{,}7977 & 0{,}7564 & 0{,}8373 \\
BTXRD & Entropy & 456 & 669 & 750 & 0{,}7498 & 0{,}7313 & 0{,}8640 \\
\end{longtable}
\endgroup


Hình \ref{fig:appendix-ood-mahalanobis-seeds} biểu diễn riêng Mahalanobis-centroid vì đây là phương pháp được nhấn mạnh trong phần kết quả chính. Khoảng biến thiên lớn hơn trên BTXRD, đặc biệt ở AUROC-OOD, cho thấy kết quả cross-dataset nhạy với trạng thái biểu diễn được học ở từng hạt giống; Semantic OOD có AUROC-OOD gần mức ngẫu nhiên và FPR@95\%TPR cao ở cả ba lần chạy.

\begin{figure}[H]
\centering
\includegraphics[width=\textwidth]{generated/appendix_seed_results/ood_mahalanobis_by_seed.png}
\caption{Các độ đo OOD của Mahalanobis-centroid tại từng hạt giống trên Semantic OOD và BTXRD}
\label{fig:appendix-ood-mahalanobis-seeds}
\end{figure}
